\begin{theorem}[High stationary classes: odd non-stable range]
\label{mod:parameters-highintermediate}
Let $K/F$ be a totally ramified cyclic extension of odd prime degree
$p$, with positive lower break $t$, and put
\[
 T=t+1,\qquad r_0=\left\lfloor\frac T2\right\rfloor,
 \qquad s_0=\left\lceil\frac T2\right\rceil.
\]
Retain the explicit residue-degree-one, integral-uniformizer, and
monogenicity hypotheses used by the norm-filtration modules.  Let
$\chi_F$ be minimal in its complete norm-character orbit, let
$\chi_K=\chi_F\circ N_{K/F}$, and suppose that their actual conductors have
stationary decompositions
\[
 m=m_F(\chi_F)=2d+\varepsilon,
 \qquad m_K=m_K(\chi_K)=2d_K+\varepsilon_K.
\]
Assume the inclusive intermediate inequalities
\[
 T\le m<2T.
\]
Writing $v=m-T$, Lean proves the exact identities
\[
 m_K+(p-1)T=pm,
 \qquad m_K=T+pv,
 \qquad \varepsilon_K=\varepsilon,
\]
and, for every $\mu\in S(K/F)$ including the identity,
\[
 m_F(\mu\chi_F)=m.
\]
Thus the lower boundary $m=T$ is included, rather than silently replaced
by the strict inequality $T<m$.  The same conductor is also recorded as
the exact Herbrand successor.

Let $\psi_K=\psi_F\circ\operatorname{Tr}_{K/F}$ and choose
$\gamma_F\in F^\times$ with
\[
 v_F(\gamma_F)=m+n_F(\psi_F).
\]
Put $\gamma_K=\iota_{F,K}(\gamma_F)$.  Total ramification is retained in
the proof as the two separate facts
\[
 e(K/F)=p,\qquad f(K/F)=1.
\]
Using the multiplicative correction $-(p-1)T$ and the additive correction
$+(p-1)T$, Lean proves
\[
 v_K(\gamma_K)=m_K+n_K(\psi_K)
              =p\bigl(m+n_F(\psi_F)\bigr).
\]
Hence every adjoint below is the denominator-sensitive adjoint for the
ordered pair $(\gamma_F,\gamma_K)$.

The exact depth arithmetic gives
\[
 d\le t,\qquad r_0\le t,
 \qquad v+s_0\le d_K+\varepsilon_K,
 \qquad p m\le2(d_K+\varepsilon_K)+(p-1)T.
\]
It constructs the complete \texttt{NormPolynomialPrecision} certificate
at the actual conductors and proves that every nonterminal symmetric term
vanishes at conductor precision.  Norm representatives are selected only
at the two permitted subcritical precisions
\[
 r_0\le t\quad\hbox{and}\quad d\le t;
\]
neither $s_0$ nor $T=t+1$ is passed to the norm-representative theorem.

More precisely, let $\tau$ be the index-one generator in the proved cyclic
enumeration of $S(K/F)$.  Lean existentially and independently chooses
supplied representatives $a,b$ of the quotient classes
\[
 \operatorname{St}_{\gamma_F,s_0}(\tau)
   \in\mathfrak p_F^v/\mathfrak p_F^{v+r_0},
 \qquad
 \beta_F=\operatorname{St}_{\gamma_F,d+\varepsilon}(\chi_F)
   \in\mathcal O_F/\mathfrak p_F^d,
\]
together with $\alpha_1,\beta_1\in K^\times$.  If
\[
 \alpha_N=N_{K/F}(\alpha_1),\qquad
 \beta_N=N_{K/F}(\beta_1),
\]
then
\[
 [a]=[\alpha_N]
   \text{ in }\mathfrak p_F^v/\mathfrak p_F^{v+r_0},
 \qquad
 [b]=[\beta_N]
   \text{ in }\mathcal O_F/\mathfrak p_F^d,
\]
and
\[
 v_K(\alpha_1)=v_F(\alpha_N)=v,
 \qquad v_K(\beta_1)=v_F(\beta_N)=0.
\]
The returned witnesses are predicates on existential choices.  The
imported source-ambiguity theorem shows that any two valid choices differ
by $U_K^{r_0}$ and $U_K^d$, respectively; no choice is promoted to a
canonical field-valued stationary parameter.

The selected exact norm $\alpha_N$ gives the same linearization as $a$ on
$\mathfrak p_F^{s_0}$.  For every $y\in\mathfrak p_K^{s_0}$ Lean then
proves the manuscript's conversion in the direction $K\to F$,
\[
 \psi_F\!\left(\frac{\alpha_NN_{K/F}(y)}{\gamma_F}\right)
 =
 \psi_F\!\left(\frac{-\alpha_N\operatorname{Tr}_{K/F}(y)}
                         {\gamma_F}\right).
\]
The proof uses
\[
 v_F(N(y))=v_K(y),\qquad
 v_F(\operatorname{Tr}(y))
 \ge\left\lfloor\frac{s_0+(p-1)T}{p}\right\rfloor\ge s_0,
\]
the exact denominator $m+n_F(\psi_F)=T+v+n_F(\psi_F)$, and the critical
truncation $N(1+y)-1\equiv\operatorname{Tr}(y)+N(y)\pmod{\mathfrak p_F^T}$.

The literal integral candidate
\[
 c_K=\iota_{F,K}(\beta_N)
       -\beta_1\frac{\iota_{F,K}(\alpha_N)}{\alpha_1}
 \in\mathcal O_K
\]
satisfies the two quotient equalities
\[
 [c_K]
 =P^*_{K/F;\gamma_F,\gamma_K}([\beta_N])
 =\operatorname{St}_{\gamma_K,d_K+\varepsilon_K}(\chi_K)
 \quad\text{in }\mathcal O_K/\mathfrak p_K^{d_K}.
\]
The first equality is proved directly from the norm--trace conversion; the
second invokes stationary-class transport under norm at the actual upstairs
conductor.  In particular, the adjoint points from the downstairs
coefficient quotient to the upstairs coefficient quotient and retains both
denominators.

For $j\in\mathbf F_p$, let $[j]$ denote the Teichmuller lift used in the
manuscript.  Lean proves both $[j]^p=[j]$ and the exact quotient bridge from
this lift to the natural multiple of the index-one class.  With the actual
twist datum for $\tau^j\chi_F$, it proves
\[
 \bigl[N_{K/F}(\beta_1+[j]\alpha_1)\bigr]
 = [\beta_N+[j]\alpha_N]
 = \operatorname{St}_{\gamma_F,d+\varepsilon}(\tau^j\chi_F)
 \quad\text{in }\mathcal O_F/\mathfrak p_F^d.
\]
The identity index $j=0$ remains in the table; its power-zero evaluation
reduces to the original $\beta_N$ class without constructing a stationary
class for the trivial norm character.
At the upper endpoint, where the comparison depth can be $T$, Lean uses
the restriction of the already constructed $s_0$-class and does not
manufacture a stationary class for $\tau$ at depth $T$.

The Teichmuller congruence itself is also exposed one layer more sharply
when $p$ and $T$ are odd and $T>1$.  Under the same trace-ideal lower-bound
hypothesis used above, for every $j\in\mathbf F_p$ Lean proves
\[
 [j]-j_{\mathbf Z}\in\mathfrak p_F^{\lfloor T/2\rfloor+1}.
\]
Indeed, writing $T=2r+1$, primality and oddness give $p\ge3$, while $T>1$
gives $r\ge1$, and hence
\[
 p(r+1)\le(p-1)(2r+1).
\]
The trace bound therefore gives $v_F(p)\ge r+1$, after which the existing
arbitrary-depth Teichmuller congruence applies directly.  The original
half-depth theorem remains available unchanged.

For the later product and phase reductions, Lean also exposes the literal
normalization
\[
 u=\frac{\beta_1}{\alpha_1},\qquad n=N_{K/F}(u),\qquad
 v_K(u)=v_F(n)=-v,
\]
together with the field identities
\[
 c_K=\iota_{F,K}(\alpha_N)(\iota_{F,K}(n)-u),\qquad
 \beta_N+[j]\alpha_N=\alpha_N(n+[j]).
\]
These are identities for the already supplied existential witnesses, not
definitions of preferred stationary representatives.

At the inclusive critical boundary $m=T$, whole-orbit minimality proves
for every nonidentity norm character that the projection of the sum of the
two stationary quotient classes to the leading quotient is nonzero.  This
excludes the only possible hidden conductor cancellation.  The more general
pre-minimality interface treats a genuine drop to an actual conductor
$q'<T$ as follows: if $q'>1$, it first constructs the new stationary class
at the stationary depth belonging to $q'$ and only then compares its
restriction with the old-layer sum; if $q'\le1$, it proves that no
stationary depth exists and constructs no stationary object.

The principal theorem packages the conductor formulas, the complete
precision certificate, the independent existential choices with their
filtration ambiguity, the norm--trace identity, the upstairs adjoint and
stationary quotient identities, and all lower twist quotient identities.
Every stationary parameter in its conclusion is a lattice-quotient class.

\LeanFile{LanglandsFirstMainLemma/Parameters/HighIntermediate.lean}
\lean{LanglandsFirstMainLemma.highParameter_intermediate_conductor_relation}
\lean{LanglandsFirstMainLemma.highParameter_intermediate_twist_conductor_eq}
\lean{LanglandsFirstMainLemma.highParameter_intermediate_ramificationIndex_eq_degree}
\lean{LanglandsFirstMainLemma.highParameter_intermediate_commonDenominator}
\lean{LanglandsFirstMainLemma.highParameter_intermediate_conductor_eq}
\lean{LanglandsFirstMainLemma.highParameter_intermediate_conductor_eq_herbrand}
\lean{LanglandsFirstMainLemma.HighIntermediateDepthArithmetic}
\lean{LanglandsFirstMainLemma.highParameter_intermediate_depthArithmetic}
\lean{LanglandsFirstMainLemma.highParameter_intermediate_variableSource_le}
\lean{LanglandsFirstMainLemma.highParameter_intermediate_precision}
\lean{LanglandsFirstMainLemma.highParameter_intermediate_intermediateTerms}
\lean{LanglandsFirstMainLemma.highParameter_intermediate_twistCrossTermBound}
\lean{LanglandsFirstMainLemma.highParameter_intermediate_normCharacterConversion}
\lean{LanglandsFirstMainLemma.highParameter_intermediate_selectedAlpha_linearization}
\lean{LanglandsFirstMainLemma.highParameter_intermediate_selectedBeta_linearization}
\lean{LanglandsFirstMainLemma.highParameter_intermediate_upstairsClass_core}
\lean{LanglandsFirstMainLemma.highParameter_intermediate_normCharacterDepth}
\lean{LanglandsFirstMainLemma.highParameter_intermediate_normRepresentatives}
\lean{LanglandsFirstMainLemma.highIntermediateTeichmullerScalar}
\lean{LanglandsFirstMainLemma.highParameter_intermediate_teichmuller_mem_integer}
\lean{LanglandsFirstMainLemma.highParameter_intermediate_teichmuller_eq_zero_iff}
\lean{LanglandsFirstMainLemma.highParameter_intermediate_teichmuller_ord_eq_zero}
\lean{LanglandsFirstMainLemma.highParameter_intermediate_teichmuller_pow}
\lean{LanglandsFirstMainLemma.highParameter_intermediate_teichmuller_congruent}
\lean{LanglandsFirstMainLemma.highParameter_intermediate_teichmuller_congruent_succ_of_odd}
\lean{LanglandsFirstMainLemma.highParameter_intermediate_teichmuller_mul_class}
\lean{LanglandsFirstMainLemma.highParameter_intermediate_norm_add_teichmuller_congruent}
\lean{LanglandsFirstMainLemma.highParameter_intermediate_index_eq_power}
\lean{LanglandsFirstMainLemma.highParameter_intermediate_indexOne_ne_one}
\lean{LanglandsFirstMainLemma.highParameter_intermediate_normCharacterClass_eq_nsmul}
\lean{LanglandsFirstMainLemma.highParameter_intermediate_twistLinearClass}
\lean{LanglandsFirstMainLemma.highParameter_intermediate_twistClass}
\lean{LanglandsFirstMainLemma.highParameter_intermediate_ratioData}
\lean{LanglandsFirstMainLemma.highParameter_intermediate_critical_noCancellation}
\lean{LanglandsFirstMainLemma.highParameter_intermediate_stationaryClass_afterDrop}
\lean{LanglandsFirstMainLemma.highParameter_intermediate_noStationaryClass_of_endpoint}
\lean{LanglandsFirstMainLemma.highParameter_intermediate}
\leanok
\SourceText{Proposition prop:high-parameter-table, part (c).}
\uses{mod:parameters-stationaryclassundernorm,
mod:parameters-minimalorbitstationary,
mod:ramification-normrepresentatives,
mod:ramification-normcharacters}
\FormalizationContract
The range is exactly $t+1\le m<2(t+1)$ and includes $m=t+1$.
The only norm-representative precisions are
$r_0=(t+1)/2\le t$ and $d\le t$; no representative is selected at
$s_0$, $t+1$, or any larger depth.  Conductors are actual least conductors,
and every twist, including the identity twist, retains its exact datum.
The norm points from $K$ to $F$, the trace in the conversion points from
$K$ to $F$, and the adjoint points from the downstairs coefficient quotient
to the upstairs coefficient quotient.  Its ordered denominator pair is
$(\gamma_F,\iota_{F,K}\gamma_F)$, and the proof retains $e=p$, $f=1$, the
different depth $(p-1)(t+1)$, and both additive and multiplicative conductor
corrections.  If a product conductor drops above one, its stationary class
is reconstructed at the new actual conductor and depth before comparison;
at conductor zero or one no stationary class is formed.  All displayed
stationary equalities are quotient equalities.  Every field element is a
supplied or existential witness with explicit filtration ambiguity, never a
canonical stationary representative.
\end{theorem}
